\documentclass[varwidth, border = 3pt]{standalone}

\usepackage{main}

%%%
% Les \packs ''ifthen'' et ''pgffor'' fournissent les macros
% ''\ifthenelse'' et ''\foreach'' d'un emploi très aisé.
% Ceci va permettre, dans un contexte technique, de tester
% l'emploi des macros proposées par le package ''main''.
%%%
\usepackage{ifthen}
\usepackage{pgffor}

\begin{document}

\csvread{mes-donnees}{data.csv}

Dimensions :
\csvdim{mes-donnees}{nbrows} lignes
et
\csvdim{mes-donnees}{nbcols} colonnes.


\medskip


Contenu de la colonne 2 d'entête "\csvdata{mes-donnees}{1}{2}" :

\bgroup
  \foreach \numrow in {2,...,\csvdim{mes-donnees}{nbrows}}{%
%%%
% ''\mathstrut'' sert à avoir des boîtes de même hauteur.
%%%
    \fbox{$\mathstrut\csvdata{mes-donnees}{\numrow}{2}$}%
    %
    \ifthenelse{\numrow = \csvdim{mes-donnees}{nbrows}}%
               {}%
               { + }%
  }
\egroup

\end{document}
